% =====================================================================
% Seção 2 — Fundamentação Teórica
% =====================================================================
\section{Fundamentação Teórica}\label{sec:fund}

Três corpos de evidência sustentam a hipótese de um fator limítrofe. Examinamos
cada um, com atenção à confiabilidade da fonte (peer-reviewed versus preprint).

\subsection{Colapso de modelo sob treinamento recursivo}

O resultado mais influente é de Shumailov \textit{et al.}~\parencite
{Shumailov2024Collapse}. Os autores mostram que, quando um modelo é treinado,
gera dados e esses dados alimentam o treino de uma geração seguinte, o processo
\emph{degenera de forma sistemática}. A cada iteração, o modelo perde os modos de
baixa probabilidade da distribuição real --- a ``cauda longa'' do mundo --- e
passa a concentrar massa em padrões de alta frequência. Formalmente, o processo
é modelado como uma cadeia de Markov que converge para distribuições
\emph{delta}: o modelo ``esquece'' o mundo real e passa a imitar, de forma cada
vez mais colapsada, aquilo que ele mesmo produz com maior frequência. O resultado
é um aumento progressivo da divergência entre a distribuição aprendida e a
distribuição original, medido por perplexidade e distâncias de distribuição.

Dois estudos independentes confirmam e aprofundam o resultado. Borji~\parencite
{Borji2024CollapseNote} demonstra, por simulações de ajuste de densidade e
reamostragem repetida, que o colapso \textbf{não é um artefato experimental}, mas
um efeito estatístico fundamental e \textbf{provavelmente inevitável} em loop
fechado: o simples ato de ajustar uma distribuição e reamostrar dela
reiteradamente produz convergência para uma gaussiana degenerada e perda de
diversidade. Dohmatob \textit{et al.}~\parencite{Dohmatob2024CollapseScaling}
interpretam o colapso como uma \emph{mudança nas scaling laws}: o treinamento
sobre dados sintéticos erase o suporte de cauda longa, alterando o expoente com
que o erro escala. Estes resultados são coerentemente apelidados de
``colapso de diversidade'' na literatura
concorrente~\parencite{Alemohammad2024SelfConsuming}.

\subsection{Scaling laws e o dado como insumo informacional}

A segunda linha parte das \emph{scaling laws}. Kaplan \textit{et
al.}~\parencite{Kaplan2020Scaling} documentaram que a perda de um modelo de
linguagem segue uma lei de potência em função do número de parâmetros, do volume
de dados e da computação, com tendências sustentadas por sete ordens de grandeza.
A revisão Chinchilla~\parencite{Hoffmann2022Chinchilla} recalibrou o regime
\emph{compute-optimal}, e trabalhos subsequentes
enfatizaram o regime \emph{data-constrained}~\parencite{Muennighoff2023DataConstrained},
em que o dado --- não a computação --- é o insumo limitante.

O aspecto crucial para a nossa questão é a origem da informação. Em 2026,
Cagnetta, Raventós, Ganguli e Wyart~\parencite{Cagnetta2026ScalingTheory}
\emph{derivaram} o expoente da scaling law limitada por dados a partir de duas
propriedades mensuráveis da linguagem natural: a taxa de decaimento da entropia
condicional ($\gamma$) e as correlações entre tokens ($\beta$), com
$\alpha_D = \gamma/(2\beta)$. Isso tem uma consequência epistemológica direta:
o expoente é função das \textbf{estatísticas da fonte real}. Um modelo que
produz IA a partir de IA não pode reabastecer essa fonte --- os dados gerados por
modelo carregam \emph{menos} informação do que o original, pois já são uma
projeção comprimida dele. Em outras palavras, o ``teto'' informacional do loop é
posto pela entropia da fonte externa; a auto-produção não a relança, apenas a
dilui.

\subsection{Auto-melhoria recursiva e o gargalo de confiabilidade}

A terceira linha reúne a literatura de \emph{self-improvement}. Desde o
\emph{bootstrapping} de raciocínios (~STaR~,~\textcite{Zelikman2022STaR}) até o
pós-treino com recompensa auto-atribuída~\parencite{Yuan2024SPIN}, a ideia de que
o modelo melhora \emph{a partir de si mesmo} é um pilar da prática atual. A
introspecção recursiva (~RISE~,~\textcite{Qu2024RISE}) alcançou ganhos de até
$\sim24\%$ em raciocínio. Contudo, duas fontes recentes revelam o limite desse
programa:

\begin{itemize}
  \item \textbf{Hacking de recompensa em self-play.}
    Bailey \textit{et al.}~\parencite{SelfPlayGuidance2026} mostram que, em loops
    longos, o modelo que \emph{gera} os problemas (\emph{Conjecturer}) aprende a
    explorar a própria função de recompensa, convergindo para problemas
    artificialmente complexos que não ajudam o modelo que \emph{resolve}
    (\emph{Solver}). O sinal de treino se corrompe quando gerador e avaliador são
    o mesmo modelo.
  \item \textbf{Discovery--reliability gap.}
    O benchmark RSIBench-Data~\parencite{RSIBench2026}, que avalia agentes de
    pesquisa \emph{data-centric} para auto-melhoria recursiva, documenta que
    agentes fazem melhorias iniciais em $58{,}33\%$ dos cenários, mas
    $78{,}26\%$ das buscas terminam abaixo do pico histórico de desempenho. O
    loop gera achados, mas \emph{não sustenta progresso monotônico} --- há ruído
    de feedback que corrói a confiança do sinal de melhoria.
\end{itemize}

\subsection{Síntese da literatura}

A leitura conjunta das três linhas aponta para um consenso emergente e
convergentes: \textbf{o teto de um loop recursivo de produção de IA é posto
pelo seu sinal de avaliação e pela qualidade informacional de seus dados.} Não há
``energia livre'' estatística: ganhos sintéticos precisam ser lastreados por
informação real. Quando \emph{gerar} e \emph{avaliar} são feitos pelo mesmo
processo, o feedback perde validade e o loop degenera ou polariza. Este é o
arcabouço que a próxima seção formaliza.
